\documentclass[11pt]{article}
\usepackage{amsmath,amssymb}
\usepackage[margin=1in]{geometry}

\newcommand{\LaG}{\mathcal L_a^\Gamma}
\newcommand{\Kamb}{K_a^{\Gamma,\mathrm{amb}}}
\newcommand{\KG}{K_a^\Gamma}
\newcommand{\PL}{P_{\mathcal L_a^\Gamma}}
\newcommand{\kappaaw}{\kappa_{a,w}}

\title{Step 191: Burnol Kappa Literature Audit}
\date{}

\begin{document}
\maketitle

\section*{Verdict}

\[
\boxed{\texttt{V\_kappa\_partial\_source\_found}}.
\]

The audit located several closely related Burnol, de Branges, and Hankel sources, but did not identify a published source that directly states the inherited target
\[
\KG(\cdot,w)=\PL\Kamb(\cdot,w),
\qquad
\kappaaw(\tau)=(T_a^*\KG(\cdot,w))(1/2+i\tau).
\]
The closest related source is Burnol's \emph{Scattering, determinants, hyperfunctions in relation to \(\Gamma(1-s)/\Gamma(s)\)}, which treats finite-interval \(J_0\) Hankel Fredholm determinants and includes a resolvent-as-reproducing-kernel observation.  That is adjacent, not yet the exact completed zeta/Sonine projection required by Steps 177--179.

\section*{Inherited Target}

Step 145 records that \(L_a\) is Burnol's extended Sonine space: functions constant on \((0,a)\) whose cosine transform is again constant on \((0,a)\).  Step 153 records the evaluator chain
\[
\mathcal U_\infty(J_a^*Y^a_{w,k})
=
T_a^*\partial_{\overline w}^{\,k}\KG(\cdot,w),
\]
and the projected-kernel identity
\[
\KG(\cdot,w)=\PL\Kamb(\cdot,w).
\]
The ambient kernel is already inherited:
\[
\Kamb(s,w)=
A_\infty(s)\overline{A_\infty(w)}
\frac{s}{s-1}\overline{\frac{w}{w-1}}
\frac{a^{s+\overline w-1}}{s+\overline w-1}.
\]
The missing record is therefore not an ambient Hardy kernel.  It is the orthogonal projection \(\PL\) onto the completed Burnol/Sonine image.

Step 178 names the missing theorem as Burnol's explicit Bessel/Hankel resolvent formula for \(\PL\).  Step 179 shows that the Branch A Calkin matrix route, Hilbert--Schmidt trace route, Weyl-sequence route, and indirect \(L_{\rho,k}\) route all require the same \(\kappa\) vector.

\section*{Burnol Sources Surveyed}

\begin{enumerate}
\item Burnol, \emph{Sur certains espaces de Hilbert de fonctions entieres, lies a la transformation de Fourier et aux fonctions L de Dirichlet et de Riemann} (2001).  The abstract states that a Sonine subspace related to zeta is constructed and that the quotient contains vectors attached to nontrivial zeros.  This is relevant to \(Y^a_{\rho,k}\), but no direct \(\PL\Kamb\) formula was identified.

\item Burnol, \emph{Sur les espaces de Sonine associes par de Branges a la transformation de Fourier} (2002).  The published abstract says explicit formulas are obtained for the de Branges functions \(E(z)\) arising from Fourier Sonine spaces.  This is the strongest positive adjacent record: knowing \(E\) gives a de Branges kernel in principle.  The audit did not locate the inherited projection resolvent \(\PL\Kamb\).

\item Burnol, \emph{On Fourier and Zeta(s)} (2004).  This source supplies the broad Fourier/co-Poisson/zeta context.  The audit did not identify the exact projected-kernel formula.

\item Burnol, \emph{Two complete and minimal systems associated with the zeros of the Riemann zeta function} (2004).  The abstract says it determines in which extended Sonine spaces zeta zeros define complete or minimal systems.  This is relevant to zero evaluators and \(Y_a\), but not to the projected-kernel resolvent.

\item Burnol, \emph{Spacetime causality in the study of the Hankel transform} (2006).  The abstract states that the Hankel transform is recovered as scattering and gives a proof of the de Branges--Rovnyak support theorem.  This is adjacent support-property machinery, not the completed zeta/Sonine projection.

\item Burnol, \emph{Scattering, determinants, hyperfunctions in relation to \(\Gamma(1-s)/\Gamma(s)\)} (2006/2008).  The abstract explicitly mentions finite-interval Fredholm determinants of the \(J_0(2\sqrt{xy})\) kernel and a resolvent-as-reproducing-kernel appendix.  This is the closest candidate for a future extraction step, but the audit did not verify it as the exact \(\PL\) formula for \(\LaG\).

\item Burnol, \emph{Entrelacement de co-Poisson} (2007).  The paper's bibliography anchors the earlier Sonine and Fourier-zeta works.  It is relevant to the co-Poisson carrier but did not expose the projected-kernel formula in the accessible metadata.
\end{enumerate}

\section*{Adjacent Classical Sources}

de Branges's 1964 work on self-reciprocal functions and the 1968 book \emph{Hilbert Spaces of Entire Functions} provide the general RKHS/de Branges kernel theory.  They do not compute the projection from the ambient Hardy kernel to Burnol's completed Sonine subspace.

The de Branges--Rovnyak support-property theorem and Burnol's Hankel causality proof are structurally close to the support constraints defining Sonine spaces.  They still do not provide the target \(\PL\Kamb\) formula without an additional identification.

Classical Bessel/Hankel sources such as Watson and the Bateman/Erdelyi tables can evaluate Bessel integrals once the relevant resolvent is known.  They do not by themselves supply Burnol's completed projection.

\section*{Specific Finding}

The audit found related formulas, not the exact one.  The only candidate strong enough to pursue is Burnol's Hankel Fredholm/resolvent paper.  A future step should extract its actual resolvent formula and compare it, variable by variable, with:
\[
\LaG=A_\infty\mathcal L_a,\qquad
\KG=\PL\Kamb,\qquad
\kappaaw=T_a^*\KG(\cdot,w).
\]
Until that comparison is performed, Path 1 is not unlocked.

\section*{Path 1 Status}

\[
\boxed{\texttt{path\_1\_status = partially\_clarified}.}
\]

The exact external-content obligation remains:
\[
\text{compute or cite }\PL\Kamb
\text{ for the completed Burnol/Sonine carrier}.
\]

\end{document}

